\subsection{Tangent cohomology and the Kuranishi line}\label{app:tangent-rederivation}
\uses{fam:A.1.2, fam:A.3.1, fam:A.4.2; the cyclic and cusp charts; \cite[III.4.10, III.4.12]{BS76}; \cite[pp.~261, 264--266]{RRV71}; \cite[p.~91, Proposition 2]{RR70}.}
\begin{proof}[Re-derivation of the known result \ref{fam:A.5.2}]
\par\smallskip\noindent\textbf{Step 1.}\ This is the supplied re-derivation of the known statement
\cite[Theorems 1.3 and 6.12]{DGKXv1}; the body retains its K status.
The analytic base-change and duality applications below are supported by
the opened standard sources \cite[III.4.10, III.4.12]{BS76},
\cite[pp.~261, 264--266]{RRV71}, and \cite[p.~91, Proposition 2]{RR70}.

\par\smallskip\noindent\textbf{Step 2.}\ \emph{The base and its line bundles.}
Write $X=X_c$ and form the analytic root curve
\[
 \mathfrak B=\sqrt[3]{(B,p_1)}\mathbin{\times_B}\sqrt[4]{(B,p_2)},
 \qquad \varpi:\mathfrak B\longrightarrow B.
\]
A root object is a line bundle with section whose indicated power is the
pullback of the section cutting out $p_j$. The actual equalities
$f^*p_1=3S_1$ and $f^*p_2=4S_2$ supply these objects as
$(\mathcal O_X(S_j),1_{S_j})$, and hence supply a morphism
$\bar f:X\to\mathfrak B$ with $\varpi\bar f=f$.
The explicit weighted quotient
\[
 \mathfrak B\simeq[(\mathbb C^2\setminus\{0\})/\mathbb C^*],
 \qquad \lambda(x,y)=(\lambda^3x,\lambda^4y),
 \qquad t=\frac{y^3}{y^3-x^4},
\]
identifies the roots with $(\mathcal O(4),y)$ and
$(\mathcal O(3),x)$ and the ordinary point $p_0$ with
$y^3-x^4=0$. In particular $\varpi^*\mathcal O_B(1)=\mathcal O(12)$.
This isomorphism follows on each root chart by the displayed root
sections; away from their zeros both sides have the same ordinary coarse
coordinate, so the chart isomorphisms glue.

\par\smallskip\noindent\textbf{Step 3.}\ For completeness the analytic line-bundle computation can be made on
the coarse curve. Removing the two stabilizer characters from any line
bundle makes it descend to $B$: on each root chart, averaging a local
frame whose central value is nonzero gives an invariant frame still
nonzero at the centre, and its invariant coefficients are holomorphic
in $s_j^{m_j}$. These frames give the descended line and its transition
functions. Consequently its Picard group has
generators $H,R_1,R_2$ and relations $3R_1=H=4R_2$.
Thus it is freely generated by $R_1-R_2$, denoted $\mathcal O(1)$;
$H=12\mathcal O(1)$, $R_1=4\mathcal O(1)$ and
$R_2=3\mathcal O(1)$. The chart descent just used is exact because
invariants of a finite group over $\mathbb C$ are an exact functor.
With the clockwise root generators, a local invariant section of
$\mathcal O(k)$ has vanishing exponents
\[
 r_3(k)\in\{0,1,2\},\quad r_3(k)\equiv k\pmod3,
 \qquad
 r_4(k)\in\{0,1,2,3\},\quad r_4(k)\equiv-k\pmod4.
\]
Dividing by those local powers gives the precise coarse direct image
\begin{equation}\label{fam:eq:ab-root-degree}
 \varpi_*\mathcal O(k)=\mathcal O_B\bigl(d(k)\bigr),\qquad
 d(k)=\frac{k}{12}-\frac{r_3(k)}3-\frac{r_4(k)}4.
\end{equation}
Indeed $k-4r_3-3r_4=12d(k)$ identifies the line bundle with
$\varpi^*\mathcal O_B(d(k))\otimes R_1^{r_3}\otimes R_2^{r_4}$.
The canonical root sections generate the invariant direct image of
the last two factors; division by their local powers is the displayed
line-bundle isomorphism on the coarse charts. Finite invariants also give
$R^q\varpi_*=0$ for $q>0$, as seen on each Stein quotient chart.
For $k=0$ the degree is zero; for
$k=1,-1,-2,-5,-6$ it is $-1$. It follows directly from the
cohomology of $\mathbb P^1$ that these five line bundles are acyclic,
whereas $H^0(\mathcal O)=\mathbb C$ and $H^1(\mathcal O)=0$.
The local equation $dt_j=m_js_j^{m_j-1}ds_j$ gives
\[
 K_{\mathfrak B}=\mathcal O(-24+2\cdot4+3\cdot3)
 =\mathcal O(-7),\qquad
 T_{\mathfrak B}(-p_0)=\mathcal O(-5).
\]

\par\smallskip\noindent\textbf{Step 4.}\ \emph{The actual base changes and cusp cohomology.}
On a finite root disc $\Delta_{s_j}$ the local presentations are
$[J'_j/\mu_{m_j}]\to[\Delta_{s_j}/\mu_{m_j}]$.
The fibre-product identity
\[
 [J'_j/\mu_{m_j}]
 \mathbin{\times_{[\Delta_{s_j}/\mu_{m_j}]}}\Delta_{s_j}
 \simeq J'_j
\]
is induced by the quotient torsor and its equivariant map to the disc;
its inverse sends a torsor with a chosen root-chart trivialization to
the corresponding point of $J'_j$. The action on $J'_j$ is free by
fam:A.1.2. Thus this is an etale atlas base change and is exactly the
smooth torus family, with its given deck generator. Over the cusp the
base change is the original cusp family. Locally it is
$t_c=z_0z_1z_2$; the homomorphism
$\mathbb C\{t_c\}\to\mathbb C\{z_0,z_1,z_2\}$ is flat because its
target is torsion-free over the discrete valuation ring.
These charts prove that $\bar f$ is representable, proper and flat,
with connected Gorenstein fibres of dimension two. Properness on the
cusp is the compact representative construction of fam:A.3.1, and on the
finite charts it is the compact torus quotient. Also
$\bar f^*p_0=W$ as reduced Cartier divisors.

\par\smallskip\noindent\textbf{Step 5.}\ Let $\nu:P_6\to W$ be the normalization, $P_6=\mathrm{dP}_6$, let
$C_h,C_v,C_d$ be the three double-curve normalizations, and let
$P_+,P_-$ be the two triple points. Each $C_e$ is $\mathbb P^1$.
Write $\rho_e^0,\rho_e^1:C_e\to P_6$ for its two paired boundary maps
and $j_e=\nu\rho_e^0=\nu\rho_e^1$. The sheaf sequence is
\begin{equation}\label{fam:eq:ab-cusp-resolution}
 0\longrightarrow\mathcal O_W\longrightarrow\nu_*\mathcal O_{P_6}
 \xrightarrow{\delta^0}\bigoplus_{e=h,v,d}j_{e*}\mathcal O_{C_e}
 \xrightarrow{\delta^1}\mathbb C_{P_+}\oplus\mathbb C_{P_-}
 \longrightarrow0,
\end{equation}
where $\delta^0(s)_e=(\rho_e^1)^*s-(\rho_e^0)^*s$ and
$\delta^1$ evaluates $t_h-t_v+t_d$ at the two ordered triple points.
Its exactness is local: for $xyz=0$ it is the alternating restriction
complex of the three coordinate planes, their pairwise lines and their
common point. Decomposing convergent power series by the variables
occurring in each monomial gives the augmented simplex complex, which
is exact. The one- and two-branch cases are the same calculation with
fewer planes. This verifies exactness for the actual boundary maps,
including the torus multipliers in their gluing.

\par\smallskip\noindent\textbf{Step 6.}\ All positive-degree structure-sheaf cohomology of $P_6,C_e$ vanishes.
Consequently \eqref{fam:eq:ab-cusp-resolution} computes cohomology by
\[
 \mathbb C\xrightarrow{0}\mathbb C^3
 \xrightarrow{\left(\begin{smallmatrix}1&-1&1\\1&-1&1\end{smallmatrix}\right)}
 \mathbb C^2.
\]
It gives $(h^0,h^1,h^2)(\mathcal O_W)=(1,2,1)$.
We also need its multiplication. The branch-intersection resolution
is the augmented semisimplicial sheaf algebra with its restriction face
maps; its local exactness makes the augmentation a multiplicative
quasi-isomorphism. Its Alexander--Whitney product therefore computes
the sheaf-cohomology cup product. In the ordered triangles
$[0,e_1,e_2]$ and $[e_1,e_2,e_1+e_2]$ it is
\[
 \alpha\smile\beta=(\alpha_h\beta_d,\alpha_d\beta_h),
 \qquad [\alpha]\smile[\beta]
 =\alpha_h\beta_d-\alpha_d\beta_h
 \quad\hbox{in }\mathbb C^2/\mathbb C(1,1).
\]
The cocycles $(1,1,0)$ and $(0,1,1)$ pair to one. Thus
$\bigwedge^2H^1(\mathcal O_W)\to H^2(\mathcal O_W)$ is an
isomorphism, including its product rather than just its dimensions.

\par\smallskip\noindent\textbf{Step 7.}\ \emph{Vertical fields and the direct images.}
Set $\mathcal T=T_X(-\log W)$ and
$\mathcal V=\ker(d\bar f^{\log})$.
The vertical transitions on the core are precisely $R_g(z)$.
Finite collar translations have vertical differential the identity,
and at the cusp the frame $\partial_{\zeta_i}$ corresponds to
$2\pi i x_i\partial_{x_i}$. The cusp transformations
$x\mapsto u(t_c)t_c^{B_0\lambda}x$ fix the coefficients of these
relative logarithmic fields. These exact identifications construct
a bundle $\mathcal E_{\rm tan}$ on $\mathfrak B$ with
$\mathcal V=\bar f^*\mathcal E_{\rm tan}$.
The fixed vector $e=(0,1)^t$ is nonzero in every logarithmic frame,
so there is an exact sequence
\[
 0\longrightarrow\mathcal O\xrightarrow{e}\mathcal E_{\rm tan}
 \longrightarrow L\longrightarrow0.
\]
The factors of $L$ are $-1/\tau$ and $1/\tau$. Therefore the usual
weight-twelve modular discriminant makes $1/\Delta_{\rm mod}(\tau)$ a
meromorphic section of $L^{12}$. It has no finite zero or pole,
and its cusp pole is simple: $\Delta_{\rm mod}(\tau)=t_c\,u(t_c)$ with
$u(0)\ne0$. Hence $L^{12}=\mathcal O(-p_0)=\mathcal O(-12)$.
The Picard calculation gives $L=\mathcal O(-1)$, and the extension
splits since $H^1(\mathcal O(1))=0$. Thus
\[
 \mathcal E_{\rm tan}\simeq\mathcal O\oplus\mathcal O(-1).
\]
There is also the exact sequence
\begin{equation}\label{fam:eq:ab-log-tangent}
 0\longrightarrow\bar f^*\mathcal E_{\rm tan}\longrightarrow\mathcal T
 \longrightarrow\bar f^*\mathcal O(-5)\longrightarrow0.
\end{equation}
On the smooth root charts it is the relative tangent sequence. At the
cusp its last map sends $\sum a_i z_i\partial_{z_i}$ to
$(\sum a_i)t_c\partial_{t_c}$, so it is onto and its kernel is the
two height-zero lattice directions already used above. Taking its
determinant and using $W=\bar f^*p_0$ gives the actual relative
dualizing line
\begin{equation}\label{fam:eq:ab-relative-dualizing}
 \omega_{X/\mathfrak B}=\bar f^*(\det\mathcal E_{\rm tan})^{-1}.
\end{equation}
In particular this computation gives
$K_X=\bar f^*\mathcal O(-6)=\mathcal O_X(-2S_2)$, with the last
identification supplied by $\bar f^*R_2=\mathcal O_X(S_2)$;
it agrees through this actual pullback with fam:A.1.1.

\par\smallskip\noindent\textbf{Step 8.}\ The fibres of $\bar f$ all have structure-sheaf cohomology dimensions
$(1,2,1)$: on ordinary and finite root charts they are tori, and at
the cusp we have just computed them. Analytic proper flat base change
on these complex-space charts gives locally free $R^q\bar f_*\mathcal O_X$
with these fibres. Its canonical comparison maps agree on the atlas
double overlaps, so they descend to $\mathfrak B$. The unit gives
$R^0\bar f_*\mathcal O_X=\mathcal O$.
The source charts are smooth and the fibres are Cartier divisors,
so relative analytic duality applies in its Gorenstein form. The
relative trace pairing, compatible on the same overlaps, and
\eqref{fam:eq:ab-relative-dualizing} give
\[
 (R^2\bar f_*\mathcal O_X)^\vee
 =\bar f_*\omega_{X/\mathfrak B}
 =(\det\mathcal E_{\rm tan})^{-1},
 \qquad R^2\bar f_*\mathcal O_X=\mathcal O(-1).
\]
Here local freeness removes the higher sheaf-Ext terms in the duality
spectral sequence. The base-change application is \cite[III.4.10 and III.4.12]{BS76}; the relative trace is supplied by \cite[pp.~261, 264--266]{RRV71}, with the Gorenstein dualizing line as in \cite[p.~91, Proposition 2]{RR70}. These replace the producing run's incorrect DGKX lemma citations.

\par\smallskip\noindent\textbf{Step 9.}\ Put $G_{\mathfrak B}=R^1\bar f_*\mathcal O_X$. Cup product and the verified fibre
comparisons show $\det G_{\mathfrak B}=R^2\bar f_*\mathcal O_X=\mathcal O(-1)$;
at $W$ this uses the nonzero pairing above.
We construct the required subline of $G_{\mathfrak B}$. On an ordinary fibre,
constant lattice cocycles identify
\[
 H^1(F_t,\mathcal O_{F_t})
 =\Lambda_{\mathbb C}^*/\operatorname{rowspan}_{\mathbb C}\Pi_c.
\]
Indeed a coboundary among constant cocycles is given by an entire
function with constant period differences. Its derivatives are
periodic holomorphic functions on the compact torus, hence constant;
the function is affine linear. This gives exactly the displayed
kernel, and both sides then have dimension two. The class of
$\gamma=(1,0,0,0)$ is nonzero, since evaluation on columns
$\widehat w,\widehat\delta$ would force both coefficients of a
putative period-row combination to vanish.
The matrices satisfy $\gamma A_j=\gamma$. Translations act trivially
on $H^1(\mathcal O)$, so this holomorphic class extends on both finite
root charts and respects their gluing.

\par\smallskip\noindent\textbf{Step 10.}\ At the cusp $\gamma$ kills
$\Lambda_{\mathrm{tor}}=\mathbb Z\langle\widehat w,\widehat\delta\rangle$.
For local sections of the free toric quotient with changes
$s_i=\Psi_{\lambda_{ij}}s_j$, the constants
$\gamma(\lambda_{ij})$ therefore form a holomorphic Cech cocycle
on the entire neighbourhood. With
$B_0=\left(\begin{smallmatrix}0&1\\-1&0\end{smallmatrix}\right)$,
the edges $e_1,e_2,e_2-e_1$ have deck vectors
$(0,1),(-1,0),(-1,-1)$. The central cocycle values are consequently
$(0,-1,-1)$, which pair to one with $(1,0,-1)$.
Its value on $W$ is nonzero. On the punctured overlap the exponential
map kills precisely $\Lambda_{\mathrm{tor}}$, so its full-lattice
cocycle is exactly $\lambda\mapsto\gamma(\lambda)$, proving equality
with the preceding class. Over the Stein cusp disc, the Leray edge
map $H^1(\bar f^{-1}D_0,\mathcal O_X)\to H^0(D_0,G_{\mathfrak B})$ is an
isomorphism, since $H^1(D_0,\mathcal O)=H^2(D_0,\mathcal O)=0$.
Thus this is a holomorphic, nonvanishing extension of $[\gamma]$.
It supplies a global subline and exact sequence
\[
 0\longrightarrow\mathcal O\xrightarrow{[\gamma]}G_{\mathfrak B}
 \longrightarrow\mathcal O(-1)\longrightarrow0.
\]
Again $H^1(\mathcal O(1))=0$ splits it. We have proved
\begin{equation}\label{fam:eq:ab-direct-images}
 R^q\bar f_*\mathcal O_X=
 \begin{cases}
 \mathcal O,&q=0,\\
 \mathcal O\oplus\mathcal O(-1),&q=1,\\
 \mathcal O(-1),&q=2.
 \end{cases}
\end{equation}

\par\smallskip\noindent\textbf{Step 11.}\ \emph{From logarithmic to ordinary tangent cohomology.}
The projection formula gives, for $q=0,1,2$ respectively,
\[
 R^q\bar f_*\mathcal V:
 \quad\mathcal O\oplus\mathcal O(-1),\quad
 \mathcal O\oplus2\mathcal O(-1)\oplus\mathcal O(-2),\quad
 \mathcal O(-1)\oplus\mathcal O(-2).
\]
For $\bar f^*\mathcal O(-5)$ the three expressions are
$\mathcal O(-5)$, $\mathcal O(-5)\oplus\mathcal O(-6)$ and
$\mathcal O(-6)$. Formula \eqref{fam:eq:ab-root-degree} makes every
nontrivial summand in this list acyclic. Leray, whose base has
cohomological dimension one, therefore gives
$h^\bullet(\mathcal V)=(1,1,0,0)$ and zero cohomology for the
horizontal term. Sequence \eqref{fam:eq:ab-log-tangent} yields
$h^\bullet(\mathcal T)=(1,1,0,0)$.

\par\smallskip\noindent\textbf{Step 12.}\ The normalization map $\nu:P_6\to X$ is an immersion on each local
branch. Let $N_\nu=\operatorname{coker}(T_{P_6}\to\nu^*T_X)$.
The exact normal-derivative map is
\begin{equation}\label{fam:eq:ab-normal-quotient}
 0\longrightarrow\mathcal T\longrightarrow T_X
 \longrightarrow\nu_*N_\nu\longrightarrow0.
\end{equation}
For a local equation $z_0\cdots z_{r-1}=0$, it sends
$\sum b_i\partial_{z_i}$ to the $r$ independent normal coefficients
$b_i|_{z_i=0}$. These can be lifted independently, proving
surjectivity; the kernel is exactly the logarithmic tangent module.
This also specifies the global map under changes of branch charts.
The cusp three-form
\[
 \Omega=d\log x_1\wedge d\log x_2\wedge dt_c
\]
is invariant under $x\mapsto u(t_c)t_c^{B_0\lambda}x$ because the
extra logarithmic differentials are multiples of $dt_c$. In any
unimodular maximal cone the height-one equation is $t_c=z_0z_1z_2$,
and the determinant-one character change gives
$\Omega=\pm dz_0\wedge dz_1\wedge dz_2$.
It therefore trivializes $K_X$ on the entire cusp neighbourhood,
including every branch intersection. Taking determinants of the
immersion tangent sequence now gives the explicit line identification
$N_\nu=K_{P_6}\otimes(\nu^*K_X)^{-1}\simeq K_{P_6}$.
Since $P_6$ is the rational projective surface $\mathrm{dP}_6$,
Serre duality gives $h^\bullet(P_6,K_{P_6})=(0,0,1)$.
The cohomology sequence of \eqref{fam:eq:ab-normal-quotient} proves
$h^\bullet(X,T_X)=(1,1,1,0)$, as asserted.

\par\smallskip\noindent\textbf{Step 13.}\ \emph{The reduced deformation germ.}
fam:A.4.2 gives a nonzero Kodaira--Spencer vector at this same $c$;
the just computed $h^1=1$ makes its map an isomorphism.
Embed the actual Kuranishi germ in its one-dimensional tangent space
and write its local analytic algebra as $\mathbb C\{t\}/I$.
The existing proper family has a local classifying map, whose pullback
sends $t$ to $g(s)$ with $g'(0)\ne0$. The holomorphic inverse-function
theorem makes $\mathbb C\{t\}\to\mathbb C\{s\}$, $t\mapsto g(s)$,
an isomorphism. Every element of $I$ maps to zero, so $I=0$.
This proves smoothness and reducedness even though $h^2(T_X)=1$.
The classifying map is consequently a local isomorphism of germs;
completeness of the Kuranishi deformation then gives the claimed
miniversality and exhaustion of sufficiently small deformations.
\end{proof}
